\subsection{Contributions}

The contributions of this work are the following.

\begin{itemize}
\item \textbf{A trading framework where one strategy implementation runs in three places.} The system has a Python library holding the strategy, the models and the market abstractions, and a C\# component running inside cTrader \cite{noauthor_forex_nodate} that connects the two through shared memory. The same strategy code runs in an offline backtest over stored market data, in the backtesting tab of cTrader, and on a live account, without modification.

\item \textbf{A market database built from quote-level data.} The historical data was collected from cTrader through the same connector and stored locally. For the seven pairs studied here it holds about 1.66 billion quotes, occupying 295 GB, together with 33.0 million bars: 32.5 million at one minute, 547,835 at one hour, 23,386 daily and 1,050 monthly, all spanning January 2014 to August 2026. The engine selects the finest series available for each run, so the price path within a bar is reconstructed from the quotes themselves rather than interpolated.

\item \textbf{A backtesting engine that charges what the account charges.} The engine steps on tick data instead of bar closes. The spread is the one quoted at the moment of the trade, read from the recorded quotes rather than assumed; the commission is computed from the traded volume at the rate of a raw-spread retail account; and stop and target levels are checked against the price path inside the bar rather than at its close. This reproduces the full cost structure of a swap-free raw-spread account, which is a product retail brokers actually offer: on such an account the spread and the commission are the two charges the trader pays, and both are applied here at the rates that account pays them. To our knowledge no published application of deep reinforcement learning to currency trading combines quote-level data, a measured spread and a real commission schedule in this way; Table~\ref{Tables:BasicRelatedWork} and Table~\ref{Tables:AdvancedRelatedWork} summarise what the reviewed studies report. It also means the returns reported in Section~\ref{Section:ExperimentalResults} are not directly comparable with results obtained on bar closes with a fixed or absent spread, and are lower than those results would be for the same policy.

\item \textbf{DDPG applied to the seven major currency pairs under one protocol.} One agent was trained per pair, with the same architecture, observation design and reward. Each was evaluated between January 2015 and January 2026 under the same cost model, so the pairs can be compared with each other and the design is tested for generality rather than tuned to one instrument.

\item \textbf{Positive and risk-controlled results on all seven pairs.} Total return over the eleven-year period ranged from $+15.7\%$ on USD/JPY to $+226.6\%$ on USD/CAD, and every pair stayed profitable when the evaluation was repeated from five different starting balances. The best pair on a risk-adjusted basis was USD/CHF, with a Sharpe ratio \cite{sharpe_sharpe_1994} of $0.80$, a Sortino ratio \cite{sortino_performance_1994} of $1.55$, a Calmar ratio \cite{bacon_practical_2023} of $0.49$ and a maximum drawdown of $8.2\%$ over 220 trades.

\item \textbf{An attribution of the returns, including one negative result.} Regressing each model's yearly return on the yearly movement of its own currency pair separates the return produced by the policy from the return produced by exposure to the market. Six of the seven pairs produced a positive intercept, and three of those had a market sensitivity close to zero. The exception is USD/JPY, whose model follows the underlying pair almost exactly and adds nothing to it. That case is reported here, together with the mechanism that causes it.
\end{itemize}